\pg{349} II. Dada a hypothese da existencia real de um só ente finito, deduz-se legitimamente a existencia do Ente infinito \footnote{Ente \emph{simples} é o que não tem partes. --- Se se considera a razão ou o fundamento da \emph{simplicidade} do ente, este pode ser, como dizemos no texto, \emph{negativo}, \emph{positivo}, e \emph{abstractivo} ou \emph{precisivo}. --- Se se considera a qualidade das partes, que se excluem do ente \emph{simples}, este divide-se principalmente, e em opposição ao \emph{composto}, em \emph{physico}, \emph{metaphysico} e \emph{logico}, conforme as partes, que se excluem, são \emph{physicas} (materia e forma), ou \emph{metaphysicas} (potencia e acto, essencia e existencia), ou \emph{logicas} (genero e differença). --- Se se considera o grau da \emph{simplicidade} do ente, este, como dizemos no texto, pode ser \emph{absoluto} e \emph{relativo}, conforme exclue --- toda e qualquer especie de \emph{composição} (\textsc{Deus}, é só Elle, é \emph{absolutamente} simples), --- ou somente uma ou outra especie (a nossa alma é \emph{relativamente} simples, porque, se exclue a composição \emph{physica}, não exclue a \emph{metaphysica}).}

ARTIGO III

\emph{Ente simples e ente composto}.

72. Ente simples e ente composto. --- Ente \emph{simples} é o que não tem partes, e por isso é indivisivel. Pode ser \emph{negativo}, \emph{positivo} e \emph{abstractivo}, conforme pode não ter partes --- ou pela tenuidade e imperfeição da sua entidade, como é o ponto mathematico, --- ou pela perfeição do seu ser, como é a nossa alma, --- ou pelo effeito da abstracção, como é o ente em geral. Neste artigo occupamo-nos do \emph{simples positivo}. --- Ente \emph{composto} é o que tem partes, e por isso é divisivel. Pode ser \emph{physico}, \emph{metaphysico} e \emph{logico}, conforme as partes, são --- \emph{physicas}, como a alma, e o corpo, --- ou \emph{metaphysicas}, como a potencia e o acto, a essencia e a existencia, --- ou \emph{logicas}, como o genero e a differença.

73. \textbf{O ente simples, quando subsiste perfeitamente, é mais perfeito que o composto.} --- O ente \emph{simples}, que subsiste perfeitamente, por não constar de partes, tem todos os elementos ou graus, que constituem a sua entidade, não só unidos, mas indivisiveis; ao passo que o ente \emph{composto}, embora tenha os seus elementos unidos, comtudo não os tem indivisiveis. Ora um ente,

(1) Nem se diga que entre o finito e o infinito não existe proporção, e por isso o conhecimento d'aquelle não pode levar ao conhecimento d'este. Por quanto, se entre o finito e o infinito não existe uma proporção de \emph{entidade} e de \emph{egualdade}, existe uma proporção de \emph{dependencia} e de \emph{analogia}; pois o finito depende do infinito, como o effeito depende da causa. --- Esta falta de proporção de \emph{entidade} e de \emph{egualdade} obsta a que tenhamos um conhecimento proprio e adequado do infinito; mas não obsta a que conheçamos a existencia do mesmo infinito, e, de algum modo, as suas perfeições.
